\documentclass[11pt,a4paper]{article}

\usepackage[T1]{fontenc}
\usepackage[utf8]{inputenc}
\usepackage[english]{babel}
\usepackage{libertine}
\usepackage[libertine]{newtxmath}
\usepackage{microtype}
\usepackage[a4paper,top=2.6cm,bottom=2.6cm,left=2.9cm,right=2.9cm]{geometry}
\usepackage{graphicx}
\usepackage{booktabs}
\usepackage{tabularx}
\usepackage{enumitem}
\usepackage{xcolor}
\usepackage[most]{tcolorbox}
\usepackage[font=small,labelfont=bf,skip=6pt]{caption}
\usepackage{titlesec}
\usepackage{fancyhdr}
\usepackage[hidelinks,breaklinks]{hyperref}
\usepackage{xurl}
\usepackage[english=american]{csquotes}

\definecolor{accent}{HTML}{0072B2}
\definecolor{panel}{HTML}{F4F6F8}
\definecolor{rule}{HTML}{C8CDD2}

\hypersetup{colorlinks=true, linkcolor=accent, urlcolor=accent, citecolor=accent,
  pdftitle={Footprints and Handprints},
  pdfauthor={Frederik Berg}}

\titleformat{\section}{\large\bfseries}{\thesection}{0.7em}{}
\titleformat{\subsection}{\normalsize\bfseries}{\thesubsection}{0.7em}{}
\titlespacing*{\section}{0pt}{17pt plus 3pt minus 2pt}{6pt}
\titlespacing*{\subsection}{0pt}{12pt plus 2pt minus 2pt}{4pt}

\setlist[itemize]{leftmargin=1.2em, topsep=3pt, itemsep=3pt, parsep=0pt}
\setlength{\parskip}{6pt}
\setlength{\parindent}{0pt}
\linespread{1.03}

\pagestyle{fancy}
\fancyhf{}
\fancyhead[L]{\footnotesize\color{gray}Footprints and Handprints}
\fancyhead[R]{\footnotesize\color{gray}\thepage}
\renewcommand{\headrulewidth}{0.4pt}
\renewcommand{\headrule}{\color{rule}\hrule height 0.4pt width\headwidth}

\newtcolorbox{keybox}{
  colback=panel, colframe=rule, boxrule=0.5pt, arc=1.5pt,
  left=9pt, right=9pt, top=8pt, bottom=8pt}

\renewcommand{\topfraction}{0.92}
\renewcommand{\bottomfraction}{0.75}
\renewcommand{\textfraction}{0.08}
\renewcommand{\floatpagefraction}{0.88}
\setlength{\floatsep}{11pt plus 2pt minus 2pt}
\setlength{\textfloatsep}{13pt plus 2pt minus 3pt}
\setlength{\intextsep}{12pt plus 2pt minus 2pt}

\newcommand{\fig}[3]{%
  \begin{figure}[htbp]
    \centering
    \includegraphics[width=0.94\linewidth]{figures/#1}
    \caption{#2}
    \label{fig:#3}
  \end{figure}}

\begin{document}

\thispagestyle{empty}

\begin{center}
  {\LARGE\bfseries Footprints and Handprints\par}
  \vspace{5pt}
  {\large How carbon and ecological accounting work,\\
   and how to read the numbers they produce\par}
  \vspace{11pt}
  {\small Frederik Berg \quad·\quad July 2026}\\[5pt]
  {\footnotesize\color{gray}Written with AI research and drafting assistance;
   see the note on sources and method at the end}
\end{center}

\vspace{5pt}
\hrule height 0.8pt
\vspace{15pt}

\section*{Summary}

Ask how large one country's carbon footprint is per person and you get several
different answers, all of them correct. For Germany the published figures run
from 6.8 to 10.4 tonnes a year. They differ because they answer different
questions: what is emitted inside the border, which gases are counted, whether
goods bought from abroad are included, and whether each resident is charged a
share of roads, hospitals and public administration. A footprint figure quoted
without those four properties cannot be compared with anything.

Two measurements share the word \enquote{footprint} and are not versions of each
other. The carbon footprint is a weight of greenhouse gases. The ecological
footprint is an area of productive land and sea, and it converts carbon dioxide
into the forest area that would be needed to absorb it. That conversion turns out
to dominate: about 60 per cent of the ecological footprint is the carbon term, and
the other five components have barely moved in sixty years. The metric works as a
way to show that consumption exceeds what the planet regenerates. It does not work
for choosing between two courses of action, because the assumption inside its
largest component moves the answer more than the options being compared.

For an individual product, what gets counted depends on the method. Tracing a
supply chain step by step is precise about what it includes and has to stop
somewhere, typically before the factory building and the machines, which
understates the total by twenty to fifty per cent. Working from the whole economy
instead misses nothing, because every purchase a company ever made is recorded
somewhere, but it treats everything within an industry as identical. Both are
standard; they fail in opposite directions.

The handprint counts emissions avoided at other people rather than emissions
caused. It exists because a personal footprint cannot show the largest thing an
individual does — persuading a landlord, an employer or a council to change
something structural. Unlike a footprint it is deliberately not exclusive: two
people can claim the same avoided tonne, so the sum of everyone's handprints can
exceed the emissions that exist. That makes it useful for describing influence and
invalid to subtract from a footprint.

\vspace{4pt}
\hrule height 0.4pt

\section{Two measurements, one word}

The \textbf{carbon footprint} is a weight. It counts greenhouse gases and reports
them in tonnes of carbon dioxide equivalent, written t\,CO$_2$e. Methane and
nitrous oxide are converted into the amount of CO$_2$ that would cause the same
warming over a hundred years, so everything lands in one unit. Climate targets are
written in this unit, which is why it is the one that matters when a decision
depends on the answer. The rules for measuring a product this way are set by
ISO~14067, a narrowing of the general life-cycle assessment standards to the single
question of climate impact \cite{iso14067,iso14040,iso14044}.

The \textbf{ecological footprint} is an area. It asks how much biologically
productive land and sea a population needs — to grow its food, catch its fish, cut
its timber, stand its buildings on, and absorb its carbon dioxide — and reports the
answer in global hectares, a hectare of world-average productivity. That demand is
compared against \emph{biocapacity}, the productive surface actually available. The
accounts, developed by Mathis Wackernagel and William Rees and now maintained by
the Global Footprint Network with York University, track six categories: cropland,
grazing land, forest products, fishing grounds, built-up land, and the forest area
needed to absorb emissions \cite{gfn-data}.

Neither number contains the other, and neither converts into the other. The
ecological footprint does have a carbon term inside it, but that term is not the
emissions — it is a hypothetical forest that would have to exist to absorb them.
The carbon footprint, in turn, says nothing about cropland or fisheries. The two
figures for one country in one year can move in opposite directions.

\begin{keybox}
A carbon footprint asks how much something warmed the planet. An ecological
footprint asks how much of the planet's productive surface it occupied. Only the
first is denominated in the quantity climate targets are set in.
\end{keybox}

\section{Why one country has four different numbers}

Germany is a useful case because all four variants are published, and the gap
between them is larger than most of the differences people argue about.

The narrowest counts only what is physically emitted inside the border, and only
carbon dioxide: 6.77 tonnes per person in 2024 \cite{owid-co2}. Add the other
greenhouse gases — agricultural methane, nitrous oxide from fertiliser,
refrigerant gases — and the same territorial basis gives 649 million tonnes for
2025, which across 83.5 million residents is 7.8 tonnes each
\cite{uba-thg,destatis-pop}. This is the number national climate law is written
against; Germany's legal targets are at least 65 per cent below 1990 by 2030 and
net zero by 2045 \cite{uba-ziele}.

Now correct for trade. Subtract the emissions caused by things Germany makes and
sells abroad, add the emissions caused abroad by things Germany buys, and the
figure rises to 9.09 tonnes of CO$_2$ per person for 2023, against 7.02 tonnes
territorial in the same year \cite{owid-co2,owid-cons}. Roughly three tonnes of
every ten a German resident is responsible for are released outside Germany.

Finally the figure consumer calculators show: 10.4 tonnes. It is trade-corrected,
counts all greenhouse gases, adds the extra warming aircraft cause at altitude,
and gives every resident an equal share of emissions from public administration,
health, education, water, waste and civil engineering \cite{greenpeace,uba-rechner}.

\fig{f2-four-numbers}{The same country, four accounting rules, no measurement
error. The note under each bar states what that rule includes.}{four}

\subsection{A share is not a footprint}

Dividing a country's total emissions by its population produces a share, and a
share is a different object from a footprint. It includes emissions from factories
producing goods for export that no resident ever touches, and excludes emissions
from a factory abroad producing something a resident bought last week.

It also hands everyone the same value, which is wrong in a way worth quantifying:
in Germany the lowest-emitting tenth of the population sits near 7 tonnes and the
highest-emitting tenth near 17.7, more than half again above the national average
\cite{uba-frage}.

\fig{f3-trade-gap}{Germany's emissions inside its border have fallen faster than
the emissions caused by German consumption. The shaded band is the part released
abroad. Data: Our World in Data and the Global Carbon Project \cite{owid-co2,owid-cons}.}{tradegap}

\section{The same question, worldwide}

Global fossil carbon dioxide emissions reached a record 38.1 billion tonnes in
2025, with a further 4.1 billion from deforestation and other land-use change
\cite{gcb2025,gcb-essd}. Divided by population that is 4.73 tonnes of CO$_2$ per
person, or 6.67 tonnes once all greenhouse gases are counted
\cite{owid-co2,owid-ghg}.

The spread between countries is about a factor of twenty. Qatar emits 41.3 tonnes
per person, the United States 14.2, China 8.7, Germany 6.8, France 4.0 and India
2.2 \cite{owid-co2}. France sits at little more than half the German figure while
being a comparably wealthy country, almost entirely because of how it generates
electricity — a reminder of how much of a national footprint is decided by
infrastructure rather than by behaviour.

\fig{f4-countries}{Emissions produced inside a country against emissions caused by
its consumption, 2023. Rich countries import more embodied emissions than they
export; manufacturing economies do the reverse. Data: \cite{owid-co2,owid-cons}.}{countries}

Country averages, though, hide a spread that is now larger than the spread between
countries. In 2019 the poorest half of humanity accounted for 12 per cent of global
emissions and the richest tenth for 48 per cent. Since 1990 the poorest half has
been responsible for 16 per cent of the growth in emissions and the richest one per
cent for 23 per cent. Most striking for how footprints are usually measured: the
bulk of the top percentile's emissions comes from what they own and invest in
rather than from what they consume \cite{chancel}.

\fig{f5-inequality}{Shares of global greenhouse gas emissions by global income
group, 2019. The multiplier is each group's share of emissions divided by its share
of people. Data: Chancel (2022) \cite{chancel}.}{inequality}

That last point has a direct consequence. A calculator that asks about flights,
heating and diet measures consumption. For most people consumption is the dominant
term and the answer is close to right. For the top of the distribution it is not,
and the group with the largest footprint is the one such a calculator understates
most.

\section{What sits inside a personal footprint}

The 10.4-tonne figure is the trade-corrected one — the fourth bar in
Figure~\ref{fig:four}, the number a consumer calculator reports. It is worth being
explicit about how it is built, because two quite different procedures are at work
inside it and neither is the one people usually assume.

For the parts a person can report about themselves, the calculation runs
\textbf{bottom-up}: kilometres driven, litres of heating oil, kilowatt-hours of
electricity, flights taken, each multiplied by an emission factor. For the parts
nobody can report — the emissions behind a shirt, a restaurant meal, a hospital —
it runs \textbf{top-down}: national totals for that sector, divided by the
population, then scaled up or down by what the person says about their habits. Food,
general consumption and public infrastructure are computed this way, as are the base
loads inside housing and mobility \cite{uba-rechner}. Section~5 explains where those
national totals come from.

So the answer to \enquote{is this per-product or is it the country divided by
people?} is: both, in different categories. Nobody adds up the individual products
they bought, because the data to do that does not exist. What varies with a person's
answers is how far above or below the national average each category lands.

The six categories the calculator uses are where most of the interpretation happens,
and few of them are where intuition puts them.

\fig{f1-breakdown}{The average German footprint by category. The notes state what
each one contains. Sources: \cite{greenpeace,uba-rechner}.}{breakdown}

The largest single category is not heating and not driving. It is everything else
bought — clothes, electronics, furniture, services — together with the machines and
factory buildings that produced all of it. Housing comes second and means heating
and hot water, plus renting and renovation work. Household electricity is a
separate and surprisingly small category at half a tonne, which is the honest
reason a green tariff moves a personal total less than people expect. Published
savings figures for switching to renewable electricity are larger than this entire
category, for a reason Section~7.1 takes up.

Mobility counts private travel only. A train taken to a client meeting is not in
it; that emission is treated as business activity and reappears inside the
consumption category. Food includes a proportional share of restaurants and
canteens and excludes pet food, which also lands in consumption \cite{uba-rechner}.

The last category is the one no user input can change. Public infrastructure, 1.2
tonnes, is an equal share of the state itself: administration, health, education,
water supply, waste disposal, and the construction of roads, railways, bridges and
tunnels \cite{uba-rechner}. It sits underneath every resident regardless of how
they live, and it is larger than the entire long-run target discussed in
Section~7.

\section{How a footprint is actually calculated}

Two methods produce these numbers, and knowing which one produced a figure explains
most of the disagreements between published sources.

\subsection{Following the supply chain, and where it stops}

The first method traces a product backwards step by step: cotton, spinning, dyeing,
sewing, shipping. It is precise about what it includes, and it has to stop
somewhere. Raw materials, process energy, direct emissions, packaging and transport
are normally inside. The factory building, the machines, the administrative
overhead and staff commuting are normally outside — commuting because it belongs to
the organisation rather than to the product, the rest because allocating a
thirty-year building across millions of units requires assumptions that are hard to
defend.

The cost of stopping is measurable. Comparing traced inventories against
whole-economy ones showed that this method systematically underestimates, in some
product systems by half or more \cite{lenzen2000}, and a group of the field's
practitioners argued that where a study draws its boundary is not decided on
scientific grounds at all \cite{suh2004}. The commonly cited range is a 20 to 50
per cent shortfall, and it is a shortfall rather than noise: the error only ever
points one way.

Standards handle this differently, which is why two figures for one product can
legitimately disagree. PAS~2050 leaves capital goods out by default
\cite{pas2050}; the European Product Environmental Footprint method requires them
\cite{pef}. For company reporting the Greenhouse Gas Protocol requires them and is
strict about when: \enquote{companies should not depreciate, discount, or amortize
the emissions from the production of capital goods over time} — a machine's full
manufacturing emissions are booked in the year it is bought, not spread across its
working life \cite{ghgp-scope3}.

\subsection{Working from the whole economy instead}

The second method starts from national accounts — the tables recording which
industries buy what from which other industries — and traces emissions through all
of them. Nothing is cut off, because every purchase any company ever made is
recorded somewhere. This is what stands behind the German calculator, which builds
its sector averages from a model of 182 product groups \cite{uba-rechner}.

The published worked example shows what that captures. German demand for research
and development in 2020 came to 110.5 billion euros and carried 8.75 million tonnes
of CO$_2$e, an intensity of 79 kilograms per thousand euros spent. Conferences those
institutions hosted and business trips their staff took are inside that figure, as
inputs bought from the hospitality and transport sectors \cite{uba-rechner}. Nothing
has to be added by hand; the supply chain is already in the matrix.

The price of that completeness is bluntness. The same report notes that research
needing a laboratory and research needing a desk carry the same coefficient
\cite{uba-rechner}. Within a category everything is treated as identical.

Three further properties of this particular model bound what its numbers can
support. Imported goods are assumed to have been made as cleanly as domestic ones,
which understates exactly those sectors that have moved abroad. Emissions from land
use and deforestation are excluded — as are, in exchange, the sinks from
reforestation. And the data runs two to three years behind the year on the label
\cite{uba-rechner}.

\subsection{The factory question}

Whether the factory that made a product is counted therefore has a specific answer
that depends on the method, and in the whole-economy method the answer is yes with
a qualification.

German investment in machines, equipment, vehicles and non-residential buildings
came to roughly 153 million tonnes of CO$_2$e in 2021. Civil engineering is assigned
to public infrastructure, residential construction to housing, and the majority
stays inside general consumption on the reasoning that machinery and industrial
buildings exist in order to produce consumer goods. Splitting it further onto
individual products is stated to be impossible with the data that exists
\cite{uba-rechner}.

So the factory is in the number, but it is spread across the whole shopping basket
rather than attached to the item bought.

\subsection{Green electricity}

The claim that renewable electricity makes an electricity footprint zero is a
boundary question rather than a factual one. Company reporting rules require both
figures to be published: one using the average emissions of the physical grid, and
one derived from what was actually contracted, such as a supplier tariff or a
renewable certificate \cite{ghgp-scope2}.

The contracted figure can legitimately be close to zero. What it does not establish
is that the purchase caused additional renewable generation to exist — the rules
allocate existing generation to whoever contracted it and do not require the
certificate to be additional. Two customers can therefore both lower their reported
figure while total grid emissions are unchanged.

The German calculator handles the closest analogous case more conservatively.
Electricity a household generates and feeds into the grid is credited as an
avoidance \emph{at other people}, calculated against the fossil generation it
displaces, and is reported on a separate line rather than subtracted from the
household's own footprint \cite{uba-rechner}.

\section{The ecological footprint, and what it is made of}

Each component of the ecological footprint is consumption divided by a world-average
yield, converted into global hectares. Biocapacity is the matching supply figure:
area, adjusted for how productive that country's land of that type is
\cite{gfn-data}.

The carbon component works differently from the other five. Emissions are reduced by
the share the oceans absorb, and what remains is divided by the rate at which
world-average forest absorbs carbon, producing an area of forest that would be
needed but does not exist.

That component dominates the result. For Germany in 2023 — the last measured year in
the current accounts — the footprint of consumption was 4.03 global hectares per
person against a domestic biocapacity of 1.59. Carbon accounted for 2.43 of those
4.03, about 60 per cent. Cropland contributed 0.79, forest products 0.39, built-up
land 0.19, grazing land 0.17 and fishing grounds 0.05 \cite{gfn-data}.

\fig{f6-ef-germany}{Germany's ecological footprint by component since 1961, against
its own biocapacity. Almost all of the movement is the carbon layer at the bottom.
Data: Global Footprint Network and York University \cite{gfn-data}.}{efde}

Remove the carbon layer and what remains — 1.60 global hectares — is close to what
Germany's own land and sea regenerate. The five non-carbon components together have
barely moved in six decades.

Worldwide the same accounts give 2.51 global hectares of demand per person for 2025
against 1.48 available, a ratio of 1.70 that is the origin of the \enquote{1.7
Earths} figure. Demand first exceeded supply around 1970 and has not returned since
\cite{gfn-data}.

\fig{f7-earths}{World ecological footprint and biocapacity per person. Biocapacity
falls steadily because population grows faster than productive area. Data:
\cite{gfn-data}.}{earths}

Earth Overshoot Day — the date on which cumulative demand equals the year's
biocapacity — falls on 30 July 2026 \cite{eod-past}. Germany's own date, meaning
the date the world would reach if everyone consumed like a German resident, is 10
May \cite{eod-country}.

\fig{f8-overshoot}{Earth Overshoot Day since 1972, recalculated on one consistent
method. Data: \cite{eod-past}.}{overshoot}

Two properties of these dates are easy to misread. The whole series is recalculated
every year, so a date published in an earlier year will not match the current one.
When the 2026 accounts revised the ocean's carbon uptake upward, the date moved
eight days later, while genuine deterioration moved it two days earlier — a net six
days later than the previous estimate, in a year the accounts themselves describe
as the highest overshoot ever recorded \cite{eod-past}. A later date can mean a
better model rather than a better world. For the same reason, per-person hectare
figures quoted from different editions will not agree, and a comparison across
sources has to check which edition each number came from.

\subsection{The case against the metric, and what survives it}

The ecological footprint has been criticised in the peer-reviewed literature for two
decades, and the central criticism is arithmetic rather than rhetorical.

\textbf{The headline claim is narrower than it sounds.} \enquote{Humanity uses 1.7
Earths} invites the reading that we are consuming farmland, forests and fisheries at
1.7 times the rate they regrow. Break the number apart and that is not what it says.
Blomqvist and colleagues showed that none of the five non-carbon categories runs a
substantial deficit — cropland demand roughly matches cropland supply, fishing roughly
matches fish, and so on. Essentially all of the measured overshoot is the carbon
component \cite{blomqvist}. The real claim is \enquote{we emit more CO$_2$ than
forests absorb}, which is true, important, and already measurable in tonnes.

\textbf{The hectares in that component do not exist, and how many there are is a
choice.} The carbon term is the forest area that \emph{would be} needed to absorb the
emissions. No such forest is planted; it is a unit conversion. How large it comes out
depends on an assumed rate at which a hectare of forest takes up carbon, and on how
much the oceans are assumed to absorb first. Neither is a measurement of what
happened. Giampietro and Saltelli made this the core of their objection: the result
swings with assumptions that cannot be pinned down, while the supply figure it is
compared against is built with the same machinery \cite{giampietro}. The 2026 revision
described above is the effect in practice — an updated ocean figure moved Earth
Overshoot Day more than reality did.

\textbf{The conversion licenses a conclusion nobody intends.} Because emissions enter
as \enquote{forest area required}, the arithmetic says global overshoot could be
closed by planting enough fast-growing trees, with emissions unchanged. The atmosphere
does not work that way; it responds to tonnes. A metric that can be satisfied without
the physical problem improving is doing something odd with its unit.

\textbf{Adding unlike things together loses the information.} Cropland, fish and
carbon are summed into one figure in hectares, so a rise in one can mask a fall in
another and the total points at no particular action. Van den Bergh and Grazi argued
on this basis that the aggregate offers nothing usable for policy \cite{vandenbergh}.

Blomqvist and colleagues concluded that the metric as constructed is \enquote{so
misleading as to preclude their use in any serious science or policy context}
\cite{blomqvist}.

Rees and Wackernagel replied that the metric answers one specific question — how much
biocapacity does a population use against how much is available — and is being judged
against a standard it never claimed to meet. The carbon share, they argued, reflects
the real dominance of fossil fuels rather than a modelling artefact \cite{rees}. The
Global Footprint Network's own limitations note concedes that the metric does not
measure pollution, biodiversity or water, and works best inside a set of indicators
rather than alone \cite{gfn-limits}.

Both positions are correct about different things, which points to a division of
labour rather than a winner. To show a general audience in one number that
consumption exceeds what the planet regenerates, the area framing and Overshoot Day
do work that a tonnage cannot. To choose between two courses of action, the
conversion to area adds an assumption without adding information, and the
sensitivity of that assumption is larger than the differences being weighed.

\section{What one person can change, and what they cannot}

Two systematic reviews have measured the size of individual actions. A synthesis of
53 studies covering 771 options found the largest ones in transport — living without
a car, switching to an electric vehicle, dropping one long return flight, each above
1.7 tonnes a year — followed by renewable electricity at 1.6 tonnes, home
refurbishment at 0.9 and a vegan diet at 0.8 \cite{ivanova}. An earlier review of 148
scenarios put living car-free at 2.4 tonnes a year, avoiding one transatlantic return
flight at 1.6 and a plant-based diet at 0.8 — four times more than comprehensive
recycling and eight times more than changing lightbulbs \cite{wynes}.

\fig{f9-levers}{What each of the largest private actions is worth per year. Grey
marks the two actions public campaigns push hardest. Sources: \cite{ivanova,wynes}.}{levers}

\subsection{The ranking is a population ranking, not a personal one}

Every figure in that chart is a median across studies, measured against an
average-emitting baseline. That makes the ordering correct for a population and
unreliable for any particular person, because \textbf{the saving an action offers is
capped by what that person already emits in the category.} Living car-free is worth
2.4 tonnes to someone who drives an average distance and nothing at all to someone
who has no car. Ivanova and co-authors treat this explicitly, modelling geographical,
technical and socio-demographic context as a determinant of how much any option
delivers \cite{ivanova}.

The chart in front of this section contains its own example. Switching to renewable
electricity appears there at 1.6 tonnes — a median across studies from countries where
household electricity is a much larger share of the footprint than it is in Germany,
and where in several cases electricity also supplies the heating. A German resident's
entire household electricity category is 0.5 tonnes (Section~4). A saving cannot
exceed the category it comes out of, so 0.5 is the ceiling on what that switch can
deliver there, whatever the median says.

The effect is large enough to reorder the list. For a person without a car the
chart's biggest item is worth nothing at all, which moves it from first place to
last. Diet shows the same dependence in a milder form: measured against one
consistent set of diets, the move from a medium meat-eater to a vegan is worth about
1.7 tonnes a year, while the move from a vegetarian to a vegan is worth about 0.6 —
7.04, 4.16 and 2.47\,kg CO$_2$e per 2000\,kcal respectively, annualised
\cite{scarborough}. Same action, two starting points, a factor of nearly three
between them.

For a German reader the diet step can be sized directly, and against something people
have an intuition for. The calculator puts the average German food footprint at 1.6
tonnes a year (Section~4), and the environment agency estimates that moving from the
average German diet to a vegan one cuts diet-related emissions by 38 to 52 per cent
\cite{uba-tierhaltung}. Combining the two gives roughly 0.6 to 0.8 tonnes a year. A
return flight inside Europe of about 1000 kilometres each way — Frankfurt to
Barcelona, Berlin to Rome — is some 2000 passenger-kilometres, and German aviation
data puts short-haul economy travel at 195 grams of CO$_2$e per passenger-kilometre
once non-CO$_2$ effects and fuel supply are counted \cite{ifeu}, so about 0.4 tonnes.
Giving up animal products for a year is therefore worth somewhere between one and a
half and two such flights. Neither side of that comparison includes emissions from
land-use change, which the calculator omits throughout and which is the larger
omission on the food side \cite{uba-rechner}.

So the published ordering describes a population, and which lever is largest for a
given person depends on where that person already stands. A generic ranking names the
categories that are usually large; only a personal baseline shows which of them still
has room left in it.

\subsection{Why the widely quoted child figure is missing from that chart}

The same review lists a fifth action far above the others — having one fewer child, at
58.6 tonnes a year \cite{wynes}. It is absent from the chart because it is a different
kind of quantity from the rest, not simply a larger one.

The study underneath it computes a lifetime carbon legacy: under current conditions in
the United States, each child adds about 9441 tonnes of CO$_2$ to an average woman's
legacy, which is 5.7 times her own lifetime emissions \cite{murtaugh}. Of that, half
of a child's emissions are assigned to each parent, a quarter of the grandchildren's,
and so on down the family tree \cite{wynes}. Three things follow. It is a lifetime sum
presented on an annual axis, where every other figure here is a saving that recurs
each year. The child is charged its own emissions again once grown, so the same tonne
sits in two ledgers, where every other figure here belongs to one person only. And it
holds emission rates at today's level for as long as the family tree runs, which is
the opposite of what the targets the literature exists to serve assume. None of this
touches the four findings above, which are ordinary annual savings measured the
ordinary way.

Adding the realistic levers together, allowing for the overlap between them and for
the ceilings just described, gets somewhere around five to seven tonnes off a
10.4-tonne starting point. The total is indicative rather than exact — it is assembled
from cross-country medians, and the German ceiling on one of its terms alone moves it
by a tonne. The conclusion survives that imprecision comfortably: it is a substantial
change in how a person lives, and it still does not reach the target.

The target itself is not a single published constant. Germany's environment agency
states that no exact value can be derived, because it depends on population growth
and on how many carbon sinks can be built, but that it must lie well below one tonne
per person per year for net zero to work — and that reaching it requires cutting more
than 90 per cent from today's level \cite{uba-frage}. An independent check gives the
same order of magnitude: 170 billion tonnes of CO$_2$ remain in the budget for 1.5
degrees \cite{gcb2025}, which across roughly 8.2 billion people is about 21 tonnes
each in total — under five years at the current German rate.

What is left after the private levers are exhausted is the electricity grid, the
building stock, the railways, the freight system and industrial process heat. None of
those is a consumption decision. The 1.2-tonne public share is not reachable by any
of them either.

\section{The handprint}

The handprint answers the question that follows directly: if most of the remaining
gap is structural, what is one person supposed to do about it?

It counts what someone changes at other people. Persuading a landlord to insulate a
building cuts emissions for every household in it. Getting an employer to put solar
on the roof, or a council to change a bus route, moves far more than any private
decision. None of it appears in a personal footprint, which by construction only
counts what passes through one person's own life. The concept originated at the
Centre for Environment Education in India and was developed in Germany by Germanwatch
with Brot für die Welt, who frame it as helping people \enquote{shape social and
political change in their environment and thus change structural conditions in the
long term} \cite{germanwatch,handabdruck}. Its corporate equivalent, avoided
emissions, is defined as the difference a low-carbon product makes against a scenario
without it, and is explicitly positioned as sitting outside a company's own emissions
inventory \cite{wbcsd}.

The property that makes the concept work is also the one that limits it, and it is
easiest to see with numbers.

Take a block of ten flats. One tenant spends a year pushing the landlord into
insulating it, and afterwards the building burns half a tonne less CO$_2$e per flat
per year. Five tonnes are genuinely avoided.

On the footprint side the bookkeeping is clean. Each of the ten residents now heats
with less gas, so each of their personal footprints falls by 0.5 tonnes. Ten times
0.5 is 5, which is exactly what was avoided. Nothing is missing and nothing is
counted twice — that is what a footprint is: every real tonne assigned to exactly one
person, totals matching reality.

On the handprint side the tenant who organised it claims 5 tonnes, and that claim is
honest too. Without them the insulation does not happen and none of the five tonnes
is avoided.

Now add the two together and you get 10 tonnes from an event that saved 5. The
tenant's own footprint also fell by 0.5, and that half-tonne now appears both in
their footprint and inside their handprint. None of this is an error. The two numbers
answer different questions — \emph{how much do I cause?} and \emph{how much did I
change?} — and the second question cannot be answered without letting several people
take credit for the same tonne.

\begin{keybox}
The footprint concept assigns all of a population's emissions to individuals without
overlap. The handprint cannot: avoidance achieved at other people cannot be assigned
unambiguously to one actor and may be claimed by several at once, so the sum of all
handprints can exceed the emissions that actually exist. Double counting is excluded
from the footprint by construction and is inherent to the handprint — which is not a
problem as long as the two are reported separately \cite{uba-rechner}.
\end{keybox}

The rule that follows is hard: a handprint may be reported beside a footprint and
never subtracted from it. The German calculator is built accordingly and quantifies
only three cases — electricity fed into the grid, donations to climate projects, and
climate-directed investments. A compensated flight still appears at full weight under
mobility, with the compensation shown separately and negatively \cite{uba-rechner}.

Once that split is clear, a large class of public claims sorts itself out.
\enquote{Our product saved four million tonnes} is a handprint claim and says nothing
about what the company emits. \enquote{We cut our emissions by four million tonnes}
is a footprint claim and can be checked against a balance that has to add up. Claims
that slide between the two are the standard shape of greenwashing, which is why the
corporate guidance treats separate reporting as a condition of credibility
\cite{wbcsd}.

\section{Where the idea came from}

A recurring argument holds that the personal carbon footprint is a public relations
construct and should be distrusted on that basis. The documented facts support part
of it. BP hired the advertising agency Ogilvy \& Mather, ran an award-winning
campaign from 2000, and launched a personal carbon footprint calculator in 2004
aimed at showing individuals how their daily life drives warming \cite{mashable}.

Two corrections matter. The term predates the campaign: the earliest use recorded by
the Oxford English Dictionary is a 1999 issue of a BBC vegetarian cooking magazine,
and the term became Oxford's word of the year in 2007 \cite{euronews}. The accounting
predates it by longer — the ecological footprint dates from the early 1990s
\cite{gfn-data}, and trade-corrected emission accounting is a standard application of
input-output analysis, a technique from national statistics.

Where a word came from is not evidence about whether a measurement works. A
statistical office computing consumption-based emissions performs the same arithmetic
regardless of who once found the vocabulary convenient. What the objection does get
right is the framing effect: a personal footprint presented without context invites
the conclusion that the sum of private choices is the solution, and Section~7 puts
that at roughly half. The response that follows is to keep the measurement and fix
the framing, which is what the handprint was constructed to do.

\section{Which metric answers which question}

\begin{table}[htbp]
\centering
\small
\caption{Every row is a legitimate metric. The failure mode is using one to answer
another's question.}
\vspace{2pt}
\begin{tabularx}{\linewidth}{@{}p{3.4cm}p{2.3cm}X@{}}
\toprule
\textbf{Question} & \textbf{Metric} & \textbf{Why, and what it cannot do} \\
\midrule
Is a country meeting its legal target? & Territorial, t\,CO$_2$e & The basis of national inventories and climate law. Ignores trade entirely. \\
\addlinespace[2pt]
How much warming does a population's consumption cause? & Consumption-based, t\,CO$_2$e & Nets out exports, adds imports. Runs two to three years behind and inherits the bluntness of whole-economy methods. \\
\addlinespace[2pt]
Which of two products is better? & Product footprint, t\,CO$_2$e & Comparable only if boundary, standard and method match. Traced figures are biased low. \\
\addlinespace[2pt]
Where should one household act first? & Consumer calculator, t\,CO$_2$e & Ranks private levers correctly. Contains a floor it cannot address and understates high-wealth emitters. \\
\addlinespace[2pt]
Does humanity exceed what the planet regenerates? & Ecological footprint, gha & Communicates overshoot in one figure. Dominated by its carbon term; unsuitable for ranking options. \\
\addlinespace[2pt]
What did an action change beyond the actor? & Handprint & Captures structural leverage. Never summed, never netted against a footprint. \\
\bottomrule
\end{tabularx}
\end{table}

\section{Conclusion}

A footprint figure without its accounting rule is not information. The four German
numbers span 6.8 to 10.4 tonnes and every one of them is correct. Before comparing
any two figures, three properties have to be fixed: which gases, which boundary, and
whether trade is counted. That single habit removes most of the confusion in public
reporting, and it applies equally to hectare figures quoted from different editions
of the ecological footprint accounts.

The metric should follow the question rather than the audience. Where a decision
depends on the answer, tonnes of CO$_2$e with a stated boundary is the defensible
unit. The ecological footprint earns its place as a way of making overshoot legible
and loses it as soon as it is asked to rank options, for the reason its critics
identified and its own accounts confirm: about 60 per cent of it is a carbon number
in an area unit, and the rest has barely moved in sixty years.

The footprint and the handprint are complements, and the property that makes the
handprint useful is the same one that makes it unsummable. Reported side by side they
describe both what a person consumes and what a person changes. Netted against each
other they describe nothing.

If one practical step follows from all of this: compute a personal footprint once, to
find the two or three levers that are actually large and to see the size of the floor
underneath them. That floor is where the remaining majority sits, and it is not
reachable by consumption decisions at all.

\vspace{10pt}
\hrule height 0.4pt
\vspace{4pt}
{\footnotesize\color{gray}All charts were generated from the primary data files
accompanying this document. Every source below was retrieved in July 2026. Research
and drafting were carried out with assistance from Claude (Anthropic); every source
cited was retrieved and checked against the primary document, and responsibility for
the content rests with the author.}

\begin{thebibliography}{99}
\footnotesize

\bibitem{greenpeace} Greenpeace Deutschland. \emph{Ökologischer Fußabdruck und CO$_2$-Fußabdruck -- was uns die Berechnung bringt.} \url{https://www.greenpeace.de/klimaschutz/klimakrise/deutschland-co2-fussabdruck}

\bibitem{uba-thg} Umweltbundesamt. \emph{Treibhausgas-Emissionen in Deutschland}, 2 July 2026. \url{https://www.umweltbundesamt.de/daten/klima/treibhausgas-emissionen-in-deutschland}

\bibitem{uba-ziele} Umweltbundesamt. \emph{Treibhausgasminderungsziele Deutschlands}, 30 June 2026. \url{https://www.umweltbundesamt.de/daten/klima/treibhausgasminderungsziele-deutschlands}

\bibitem{uba-frage} Umweltbundesamt. \emph{Wie hoch sind die Treibhausgasemissionen pro Person in Deutschland und wie viel wäre klimaverträglich?} \url{https://www.umweltbundesamt.de/service/uba-fragen/wie-hoch-sind-die-treibhausgasemissionen-pro-person}

\bibitem{uba-rechner} Schunkert, S., Siewert, J., Pitz, P., Paar, A., Hertle, H., Berg, F., Dittrich, M., Dingeldey, M. \emph{Der UBA-CO$_2$-Rechner für Privatpersonen -- Hintergrundinformationen.} Umweltbundesamt, TEXTE 19/2025. \url{https://www.umweltbundesamt.de/system/files/medien/11850/publikationen/19_2025_texte.pdf}

\bibitem{destatis-pop} Statistisches Bundesamt. \emph{Bevölkerung Deutschlands sinkt im Jahr 2025 um 0,1\,\%.} \url{https://www.destatis.de/DE/Presse/Pressemitteilungen/2026/01/PD26_032_124.html}

\bibitem{owid-co2} Our World in Data, from the Global Carbon Project. \emph{CO$_2$ emissions per capita.} \url{https://ourworldindata.org/grapher/co2-emissions-per-capita}

\bibitem{owid-cons} Our World in Data, from the Global Carbon Project. \emph{Per capita consumption-based CO$_2$ emissions.} \url{https://ourworldindata.org/grapher/consumption-co2-per-capita}

\bibitem{owid-ghg} Our World in Data. \emph{Per capita greenhouse gas emissions including land use.} \url{https://ourworldindata.org/grapher/per-capita-ghg-emissions}

\bibitem{gcb2025} Global Carbon Project. \emph{Fossil fuel CO$_2$ emissions hit record high in 2025}, 13 November 2025. \url{https://globalcarbonbudget.org/fossil-fuel-co2-emissions-hit-record-high-in-2025/}

\bibitem{gcb-essd} Friedlingstein, P. et al. \emph{Global Carbon Budget 2025.} Earth System Science Data 18, 3211. \url{https://essd.copernicus.org/articles/18/3211/2026/}

\bibitem{gfn-data} Global Footprint Network and York University Ecological Footprint Initiative. \emph{Footprint Data Platform}, National Footprint and Biocapacity Accounts, 2026 edition. \url{https://data.footprintnetwork.org/}

\bibitem{eod-past} Global Footprint Network. \emph{Past Earth Overshoot Days.} \url{https://www.overshootday.org/newsroom/past-earth-overshoot-days/}

\bibitem{eod-country} Global Footprint Network. \emph{Country Overshoot Days 2026.} \url{https://www.overshootday.org/newsroom/country-overshoot-days/}

\bibitem{blomqvist} Blomqvist, L., Brook, B.W., Ellis, E.C., Kareiva, P.M., Nordhaus, T., Shellenberger, M. \emph{Does the Shoe Fit? Real versus Imagined Ecological Footprints.} PLOS Biology 11(11): e1001700, 2013. \url{https://journals.plos.org/plosbiology/article?id=10.1371/journal.pbio.1001700}

\bibitem{rees} Rees, W.E., Wackernagel, M. \emph{The Shoe Fits, but the Footprint is Larger than Earth.} PLOS Biology 11(11): e1001701, 2013. \url{https://journals.plos.org/plosbiology/article?id=10.1371/journal.pbio.1001701}

\bibitem{vandenbergh} van den Bergh, J.C.J.M., Grazi, F. \emph{Ecological Footprint Policy? Land Use as an Environmental Indicator.} Journal of Industrial Ecology 18(1), 10--19, 2014. \url{https://doi.org/10.1111/jiec.12045}

\bibitem{giampietro} Giampietro, M., Saltelli, A. \emph{Footprints to nowhere.} Ecological Indicators 46, 610--621, 2014. \url{https://doi.org/10.1016/j.ecolind.2014.01.030}

\bibitem{gfn-limits} Global Footprint Network. \emph{Ecological Footprint: Limitations and Criticism.} \url{https://www.footprintnetwork.org/content/uploads/2020/12/Footprint-Limitations-and-Criticism.pdf}

\bibitem{chancel} Chancel, L. \emph{Global carbon inequality over 1990--2019.} Nature Sustainability 5, 931--938, 2022. \url{https://www.nature.com/articles/s41893-022-00955-z}

\bibitem{wynes} Wynes, S., Nicholas, K.A. \emph{The climate mitigation gap: education and government recommendations miss the most effective individual actions.} Environmental Research Letters 12(7), 074024, 2017. \url{https://doi.org/10.1088/1748-9326/aa7541}

\bibitem{ivanova} Ivanova, D. et al. \emph{Quantifying the potential for climate change mitigation of consumption options.} Environmental Research Letters 15(9), 093001, 2020. \url{https://doi.org/10.1088/1748-9326/ab8589}

\bibitem{scarborough} Scarborough, P. et al. \emph{Vegans, vegetarians, fish-eaters and meat-eaters in the UK show discrepant environmental impacts.} Nature Food 4, 565--574, 2023. \url{https://www.nature.com/articles/s43016-023-00795-w}

\bibitem{murtaugh} Murtaugh, P.A., Schlax, M.G. \emph{Reproduction and the carbon legacies of individuals.} Global Environmental Change 19(1), 14--20, 2009. \url{https://doi.org/10.1016/j.gloenvcha.2008.10.007}

\bibitem{uba-tierhaltung} Umweltbundesamt. \emph{Fragen und Antworten zu Tierhaltung und Ernährung.} \url{https://www.umweltbundesamt.de/themen/landwirtschaft/landwirtschaft-umweltfreundlich-gestalten/fragen-antworten-zu-tierhaltung-ernaehrung}

\bibitem{ifeu} Allekotte, M., Kräck, J., Kämper, C. \emph{Harmonized GHG Emission Factors for Air Travel in Germany.} ifeu paper 01/2026, Heidelberg. \url{https://doi.org/10.5281/zenodo.19208116}

\bibitem{germanwatch} Germanwatch e.V. \emph{The Handprint Concept. Changing Structures towards Sustainability}, 2024. \url{https://www.germanwatch.org/en/90288}

\bibitem{handabdruck} Germanwatch and Brot für die Welt. \emph{Handabdruck.} \url{https://www.handabdruck.eu/}

\bibitem{wbcsd} World Business Council for Sustainable Development. \emph{Guidance on Avoided Emissions v2.0}, 24 July 2025. \url{https://www.wbcsd.org/resources/guidance-on-avoided-emissions/}

\bibitem{iso14040} ISO 14040:2006. \emph{Environmental management -- Life cycle assessment -- Principles and framework.} \url{https://www.iso.org/standard/37456.html}

\bibitem{iso14044} ISO 14044:2006. \emph{Environmental management -- Life cycle assessment -- Requirements and guidelines.} \url{https://www.iso.org/standard/38498.html}

\bibitem{iso14067} ISO 14067:2018. \emph{Greenhouse gases -- Carbon footprint of products -- Requirements and guidelines for quantification.} \url{https://www.iso.org/standard/71206.html}

\bibitem{pas2050} BSI. \emph{PAS 2050:2011 -- Specification for the assessment of the life cycle greenhouse gas emissions of goods and services.} \url{https://knowledge.bsigroup.com/products/specification-for-the-assessment-of-the-life-cycle-greenhouse-gas-emissions-of-goods-and-services}

\bibitem{pef} Commission Recommendation (EU) 2021/2279 on the use of the Environmental Footprint methods. \url{https://eur-lex.europa.eu/eli/reco/2021/2279/oj}

\bibitem{ghgp-scope3} WRI and WBCSD Greenhouse Gas Protocol. \emph{Technical Guidance for Calculating Scope 3 Emissions}, Category 2: Capital Goods. \url{https://ghgprotocol.org/sites/default/files/standards/Scope3_Calculation_Guidance_0.pdf}

\bibitem{ghgp-scope2} Greenhouse Gas Protocol. \emph{Top Ten Questions about the Scope 2 Guidance.} \url{https://ghgprotocol.org/blog/top-ten-questions-about-scope-2-guidance}

\bibitem{lenzen2000} Lenzen, M. \emph{Errors in Conventional and Input-Output-based Life-Cycle Inventories.} Journal of Industrial Ecology 4(4), 127--148, 2000. \url{https://doi.org/10.1162/10881980052541981}

\bibitem{suh2004} Suh, S. et al. \emph{System Boundary Selection in Life-Cycle Inventories Using Hybrid Approaches.} Environmental Science \& Technology 38(3), 657--664, 2004. \url{https://doi.org/10.1021/es0263745}

\bibitem{mashable} Kaufman, M. \emph{The carbon footprint sham.} Mashable, 13 July 2020. \url{https://mashable.com/feature/carbon-footprint-pr-campaign-sham}

\bibitem{euronews} Euronews Green. \emph{25 years counting `carbon footprints': Experts assess the term beloved by oil firms}, 9 November 2024. \url{https://www.euronews.com/2024/11/09/25-years-counting-carbon-footprints-experts-assess-the-term-beloved-by-oil-firms}

\end{thebibliography}

\end{document}
